\begin{abstract}
eBPF lets developers run custom programs inside the Linux kernel, where a static verifier must prove each one safe before it loads. When the proof fails, the verifier rejects the program with a single message that names the failing operation. However, the message under-determines the repair. To learn what the message leaves out, we conduct an empirical study on \bench, a corpus of 235 reproduced rejections. We find that the message leaves the developer to recover all three parts of a repair: which bug occurred, which layer must change, and where the fix goes. Although the message omits all three, the log still carries the verifier's full per-instruction reasoning. We present \sys, a userspace tool that reconstructs the discarded proof from this reasoning and locates the fix site. We then feed \sys's diagnostic to a downstream repair model and test whether the model repairs the program. Across three evaluated models on the current 75-case \testset repair suite, \sys raises one-shot repair success by 10.7--21.4 percentage points, from 0.0--37.3\% with raw logs to 10.7--50.7\% with diagnostics.
\end{abstract}
